% =====================================================================================
% REVISIÓN 1: solo los «Temas tentativos del marco teórico».
% REVISIÓN 2: el Marco teórico completo (cuando lo desarrolles, sustituye o convierte los temas tentativos en subsecciones).
% Escribe tu texto FUERA de los recuadros \begin{guia} ... \end{guia}.
% =====================================================================================

\section{Marco teórico} \label{sec:marco_teorico} % audit:req=marco_teorico

\begin{guia}[Guía: Marco teórico (Revisión 2)]
\textbf{¿Para qué sirve?} Reúne los conceptos, teorías, tecnologías, métodos y normas que se necesitan para entender y
justificar tu propuesta. No es un glosario ni un resumen de Internet: cada tema se explica \emph{en función de tu proyecto}.

\textbf{Debe responder (por cada tema):}
\begin{itemize}[nosep,leftmargin=*]
  \item ¿Qué es y cómo funciona? (definición con fuente citada y, si ayuda, un esquema o ecuación).
  \item ¿Para qué se usa en mi proyecto y por qué lo elegí frente a las alternativas?
  \item ¿Qué ventajas, limitaciones o buenas prácticas debo conocer?
\end{itemize}
Además: ¿qué metodología de desarrollo o investigación usarás y por qué? ¿qué normas o estándares aplican (calidad,
seguridad, protección de datos)?

\textbf{Organización:} una \verb|\subsection| por tema, de lo general a lo particular.

\textbf{Extensión mínima:} 400 palabras de desarrollo (sin contar los temas tentativos de la Revisión 1).

\textbf{Fuentes:} libros, artículos y documentación oficial. Cada afirmación importante lleva \verb|\cite{clave}|.

\textbf{Errores frecuentes:} temas sin relación con el proyecto; copiar y pegar; fuentes que son solo URLs de blogs; no
citar; temas que luego no se usan en el desarrollo.

\textbf{Cómo se revisa:} pertinencia, profundidad suficiente, calidad de las fuentes y que cada tema se vincule con la
propuesta.
\end{guia}

% >>> REVISIÓN 2: ESCRIBE AQUÍ el marco teórico (introducción breve y una \subsection por tema) <<<


\subsection{Temas tentativos del marco teórico} \label{sec:temas_tentativos} % audit:req=temas_marco_teorico

\begin{guia}[Guía: Temas tentativos del marco teórico (Revisión 1)]
\textbf{¿Para qué sirve?} En la primera revisión no se pide el marco teórico completo, sino la lista de los temas que
vas a desarrollar. Así el director puede confirmar que vas bien encaminado.

\textbf{Debe responder:} ¿cuáles son los temas (de 5 a 10) que necesito para sustentar el proyecto? Para cada uno, en una
o dos oraciones: ¿qué es y por qué lo necesito?

\textbf{Formato:} una lista (\verb|\begin{enumerate}| o \verb|itemize|), un tema por elemento.

\textbf{Extensión mínima:} 25 palabras. En la Revisión 2 conviertes cada tema en una subsección desarrollada.

\textbf{Cómo se revisa:} que los temas cubran lo necesario para el proyecto y no incluyan relleno.
\end{guia}

\subsubsection{Modelo predictivo}

Un modelo predictivo es una herramienta estadística que infiere tendencias futuras a partir de datos históricos. Su fundamento consiste en identificar relaciones históricas entre variables y construir estructuras matemáticas o estadísticas que permitan inferir resultados desconocidos.
El módulo contendrá un modelo predictivo con el fin de proporcionar información a la población.

\subsubsection{Análisis de Fourier}

El análisis de Fourier es la herramienta matemática fundamental para el estudio de fenómenos periódicos. Su principio básico consiste en representar cualquier función periódica como una serie o integral infinita de funciones seno y coseno. Este método en el contexto actual este metodo se utiliza para representar el modelo.

\subsubsection{Función logística}

La función logística fue propuesta originalmente por el matemático belga Pierre François Verhulst en el siglo XIX para perfeccionar el modelo de crecimiento poblacional de Malthus. Su planteamiento central es que el crecimiento poblacional no puede prolongarse indefinidamente, dado que está limitado por los recursos y el espacio, por lo que el crecimiento termina por desacelerarse y estabilizarse. Esto aplica para la economía y el crecimiento del consumo de electricidad por parte de la población y las industrias

\subsubsection{Modelo híbrido Fourier-Logístico}

La combinación del análisis de Fourier con la función logística permite construir un modelo predictivo híbrido que captura simultáneamente las fluctuaciones periódicas y el crecimiento tendencial. El modelo propuesto por Bordallo Favela y colaboradores descompone el consumo eléctrico en dos componentes: la componente periódica \( EE_C \), impulsada por la variación estacional de la temperatura, y la componente de tendencia \( EE_T \), impulsada por el desarrollo económico.

\subsubsection{Tecnologías frontend React y Vite}

El sistema de visualización desarrollado en esta investigación emplea React y Vite para construir la aplicación, con el fin de soportar la presentación interactiva de los resultados predictivos y el análisis de datos.

\textbf{React} es una biblioteca de JavaScript de código abierto desarrollada por Meta (antes Facebook) para construir interfaces de usuario. 

\textbf{Vite} es una herramienta de construcción de nueva generación, cuyo nombre proviene del francés ``rápido''. La innovación central de Vite radica en aprovechar la capacidad nativa de módulos ES (ES Modules) del navegador en el entorno de desarrollo, lo que permite una \textbf{sustitución de módulos en caliente (HMR) de muy alto rendimiento}: cuando se modifica el código fuente, solo el módulo modificado necesita recargarse, en lugar de actualizar toda la página, manteniendo así una experiencia de desarrollo fluida incluso en proyectos de gran tamaño.